\documentclass[11pt,a4paper]{article}
\usepackage[utf8]{inputenc}
\usepackage[romanian]{babel}
\usepackage{geometry}
\usepackage{xltabular}
\usepackage{multirow}
\usepackage{enumitem}
\geometry{a4paper, left=2cm, right=2cm, top=2cm, bottom=2cm}

% Setări pentru a face tabelele să arate mai bine
\renewcommand{\arraystretch}{1.3}
\setlist{nosep}

\begin{document}

\begin{center}
    \textbf{\Large FIŞA DISCIPLINEI}
\end{center}

\section*{1. Date despre program}
\noindent
\begin{xltabular}{\textwidth}{|l|X|}
\hline
1.1 Instituţia de învăţământ superior & Universitatea Transilvania din Brașov \\ \hline
1.2 Facultatea & Matematică și Informatică \\ \hline
1.3 Departamentul & Matematică și Informatică \\ \hline
1.4 Domeniul de studii & Informatică \\ \hline
1.5 Ciclul de studii & Licență \\ \hline
1.6 Programul de studiu / Calificarea & Informatică \\ \hline
\end{xltabular}

\section*{2. Date despre disciplină}
\noindent
\begin{xltabular}{\textwidth}{|l|X|l|l|l|l|l|l|}
\hline
2.1 Denumirea disciplinei & \multicolumn{7}{l|}{\textbf{Fundamentele algebrice ale informaticilor}} \\ \hline
2.2 Titularul activităților de curs & \multicolumn{7}{l|}{(Completare manuală/ulterioară)} \\ \hline
2.3 Titularul activităților de seminar & \multicolumn{7}{l|}{(Completare manuală/ulterioară)} \\ \hline
2.4 Anul de studiu & I & 
2.5 Semestrul & I & 
2.6 Tipul de evaluare & C & 
2.7 Regimul disciplinei & DI \\ \hline
\end{xltabular}

\section*{3. Timpul total estimat (ore pe semestru al activităților didactice)}
\noindent
\begin{xltabular}{\textwidth}{|l|l|l|X|l|X|}
\hline
3.1 Număr de ore pe săptămână & & Din care: 3.2 curs & 28 & 3.3 seminar/laborator & 8 \\ \hline
3.4 Total ore din planul de învăţământ & 36 & Din care: 3.5 curs & 28 & 3.6 seminar/laborator & 8 \\ \hline
\multicolumn{6}{|l|}{Distribuţia fondului de timp: \textbf{5 credite}} \\ \hline
\multicolumn{5}{|l|}{Studiul după manual, suport de curs, bibliografie şi notiţe} & 0 \\ \hline
\multicolumn{5}{|l|}{Documentare suplimentară în bibliotecă, pe platforme, pe teren} &  \\ \hline
\multicolumn{5}{|l|}{Pregătire seminarii/laboratoare, teme, referate, portofolii şi eseuri} & 20 \\ \hline
\multicolumn{5}{|l|}{Tutoriat} & 8 \\ \hline
\multicolumn{5}{|l|}{Examinări} & \\ \hline
\multicolumn{5}{|l|}{Alte activități:} & \\ \hline
\multicolumn{3}{|l|}{3.7 Total ore studiu individual} & \multicolumn{3}{l|}{89} \\ \hline
\multicolumn{3}{|l|}{3.8 Total ore pe semestru} & \multicolumn{3}{l|}{125} \\ \hline
\multicolumn{3}{|l|}{3.9 Numărul de credite} & \multicolumn{3}{l|}{5} \\ \hline
\end{xltabular}

\section*{4. Precondiţii (acolo unde este cazul)}
\noindent
\begin{xltabular}{\textwidth}{|l|X|}
\hline
4.1 de curriculum & \\ \hline
4.2 de competenţe & \\ \hline
\end{xltabular}

\section*{5. Condiţii (acolo unde este cazul)}
\noindent
\begin{xltabular}{\textwidth}{|l|X|}
\hline
5.1 De desfăşurare a cursului & • \\ \hline
5.2 De desfăşurare a seminarului/laboratorului & • \\ \hline
\end{xltabular}

\section*{6. Competenţele specifice acumulate}
\noindent
\begin{xltabular}{\textwidth}{|p{3.5cm}|X|}
\hline
Competenţe profesionale & 
\vspace{-0.3cm}
\begin{itemize}[leftmargin=*]
    \item \textbf{CP1. Creeze softuri, dezvoltă prototipul pentru software}
    \begin{itemize}
            \item R.I. 1.1. Absolventul transpune o serie de cerințe într-un concept de software clar și organizat
            \item R.I. 1.2. Absolventul creează o primă versiune incompletă sau preliminară a unei aplicații software pentru a simula unele aspecte specifice ale produsului final
        \end{itemize}
    \item \textbf{CP2. Utilizează șabloane de proiectare de software, creează diagrame de proces}
    \begin{itemize}
            \item R.I. 2.1. Absolventul utilizează soluții reutilizabile, întocmește cele mai bune practici în vederea îndeplinirii activităților comune de dezvoltare TIC în dezvoltarea și proiectarea de software
            \item R.I. 2.2. Absolventul alcătuiește o diagramă care ilustrează progresul sistematic înregistrat pe parcursul unei proceduri sau al unui sistem utilizând linii de legături și un set de simboluri
        \end{itemize}
    \item \textbf{CP3. Analizează specificațiile software, definește arhitectura software, proiectează sistemul informatic}
    \begin{itemize}
            \item R.I. 3.1. Absolventul evaluează specificațiile unui produs sau sistem software care urmează să fie dezvoltat prin identificarea cerințelor funcționale și nefuncționale, a constrângerilor și a posibilelor seturi de cazuri de utilizare care ilustrează interacțiunile dintre software și utilizatorii săi
            \item R.I. 3.2. Absolventul creează și documentează structura produselor software, inclusiv componentele, cuplarea și interfețele. Asigură fezabilitatea, funcționalitatea și compatibilitatea cu platformele existente
            \item R.I. 3.3. Absolventul definește arhitectura, componentele, modulele, interfețele și datele pentru sistemele informatice integrate (hardware, software și rețea), pe baza cerințelor și a specificațiilor sistemului
        \end{itemize}
\end{itemize}
\vspace{-0.3cm}
\\ \hline
Competenţe transversale & 
\vspace{-0.3cm}
\begin{itemize}[leftmargin=*]
    \item \textbf{CT1. Aplică competențe de bază în materie de programare, operează echipamente hardware digitale}
    \begin{itemize}
            \item R.I. 1.1. Absolventul enumeră instrucțiuni simple pentru un sistem informatic în vederea rezolvării problemelor sau a îndeplinirii sarcinilor la un nivel de bază și cu orientări adecvate, dacă este necesar
            \item R.I. 1.2. Absolventul utilizează echipamente precum monitor, mouse, tastatura, dispozitive de stocare, imprimante și scanere, pentru a efectua operații precum conectarea, pornirea, oprirea, repornirea, salvarea fișierelor și alte operațiuni
        \end{itemize}
    \item \textbf{CT2. Utilizează software de comunicare și colaborează, efectuează căutări pe Internet}
    \begin{itemize}
            \item R.I. 2.1. Absolventul utilizează instrumente și tehnologii digitale simple pentru a comunica, a interacționa și a colabora cu ceilalți
            \item R.I. 2.2. Absolventul caută date, informații și conținut, prin căutări pe Internet
        \end{itemize}
    \item \textbf{CT3. Identifică probleme, soluționează probleme}
    \begin{itemize}
            \item R.I. 3.1. Absolventul identifică și detectează diverse probleme și aspecte și ia decizii cu privire la cea mai bună cale de urmat
            \item R.I. 3.2. Absolventul raportează problemele în consecvență atunci când este necesar
            \item R.I. 3.3. Absolventul găsește soluții la probleme practice, operaționale sau conceptuale într-o gamă largă de contexte
        \end{itemize}
    \item \textbf{CT4. Demonstrează angajament, gândște rapid, gândște analitic, și menține concentrarea pentru perioade lungi de timp}
    \begin{itemize}
            \item R.I. 4.1. Absolventul demonstrează disponibilitatea de a-și asuma sarcini, chiar dacă acestea sunt dificile sau incomode
            \item R.I. 4.2. Absolventul este în măsură să înțeleagă și să prelucreze cele mai importante aspecte ale datelor și conexiunile acestora în mod rapid și precis
            \item R.I. 4.3. Absolventul găsește soluții folosind logica și raționamentul pentru a identifica punctele tari și punctele slabe ale soluțiilor alternative, concluzionând sau abordărilor problemelor
            \item R.I. 4.4. Absolventul rămâne concentrat pe o perioadă lungă de timp pentru a judeca corect și a lua decizii adecvate
        \end{itemize}
    \item \textbf{CT5. Lucrează în echipă, organizează informații, obiecte și resurse, se dovedește a doriință de învățare, se adaptează la schimbare, face față stresului}
    \begin{itemize}
            \item R.I. 5.1. Absolventul lucrează cu încredere în cadrul unui grup, fiecare făcând-și partea lui în serviciul întregului
            \item R.I. 5.2. Absolventul înțelege sarcinile care îi revin și procesele aferente. Organizează informaţii, obiecte şi resurse prin metode sistematice şi în conformitate cu anumite standarde și asigură gestionarea sarcinii
            \item R.I. 5.3. Absolventul dă dovadă de o atitudine pozitivă faţă de cerinţele noi şi provocatoare care pot fi satisfăcute doar prin învăţare pe tot parcursul vieţii
            \item R.I. 5.4. Absolventul îşi schimbă atitudinea sau comportamentul pentru a se adapta modificărilor de la locul de muncă
            \item R.I. 5.5. Absolventul gestionează provocările, perturbările şi schimbările şi se redresează în urma regreselor şi a adversităților
        \end{itemize}
\end{itemize}
\vspace{-0.3cm}
\\ \hline
\end{xltabular}

\section*{7. Obiectivele disciplinei}
\noindent
\begin{xltabular}{\textwidth}{|l|X|}
\hline
7.1 Obiectivul general al disciplinei & \\ \hline
7.2 Obiectivele specifice & \\ \hline
\end{xltabular}

\section*{8. Conţinuturi}
\noindent
\begin{xltabular}{\textwidth}{|X|l|l|}
\hline
\textbf{8.1 Curs} & \textbf{Metode de predare} & \textbf{Observaţii} \\ \hline
 & & \\ \hline
 & & \\ \hline
\multicolumn{3}{|l|}{\textbf{Bibliografie}} \\ \hline
\multicolumn{3}{|l|}{1. } \\ \hline
\end{xltabular}

\noindent
\begin{xltabular}{\textwidth}{|X|l|l|}
\hline
\textbf{8.2 Seminar/laborator} & \textbf{Metode de predare} & \textbf{Observaţii} \\ \hline
 & & \\ \hline
 & & \\ \hline
\multicolumn{3}{|l|}{\textbf{Bibliografie}} \\ \hline
\multicolumn{3}{|l|}{1. } \\ \hline
\end{xltabular}

\section*{9. Coroborarea conţinuturilor disciplinei cu aşteptările...}
\noindent (Descriere despre cum se potrivește materia cu cerințele angajatorilor și asociațiilor profesionale)
\vspace{0.5cm}

\section*{10. Evaluare}
\noindent
\begin{xltabular}{\textwidth}{|l|X|X|c|}
\hline
\textbf{Tip activitate} & \textbf{10.1 Criterii de evaluare} & \textbf{10.2 Metode de evaluare} & \textbf{10.3 Pondere nota finală} \\ \hline
10.4 Curs & & Examen (C) & \% \\ \hline
10.5 Seminar/laborator & & & \% \\ \hline
\multicolumn{4}{|l|}{\textbf{10.6 Standard minim de performanţă}} \\ \hline
\multicolumn{4}{|l|}{• Obținerea notei 5 la evaluarea finală.} \\ \hline
\end{xltabular}

\vspace{1cm}
\noindent
\begin{xltabular}{\textwidth}{X X X}
\textbf{Data completării} & \textbf{Semnătura titularului de curs} & \textbf{Semnătura titularului de seminar} \\
......................... & ......................... & ......................... \\[1cm]
\textbf{Data avizării în departament} & & \textbf{Semnătura directorului de departament} \\
......................... & & ......................... \\
\end{xltabular}

\end{document}